\section{Bounded comparison on the equation background}
\label{uni-sec:bounded}

A condition on the background gives bounded uniqueness even when the positive-density region is disconnected. Ordinary comparison and an auxiliary solution whose density is enlarged on a strict sublevel set supply this criterion. The Guan--Li condition supplies ordinary comparison by making the total Monge--Amp\`ere mass independent of the potential.

Set $V=\int_X\omega^n$. We use the bounded products of \cref{prop:bounded-products} and the notation $\mathrm{MA}_\omega$.

The Guan--Li condition
\begin{equation}
   \ddc\omega=0,\qquad \ddc(\omega^2)=0
   \label{uni-eq:GL}
\end{equation}
will be imposed on the equation background $\omega$.

\Needspace{17\baselineskip}
\begin{theorem}\label{uni-thm:GL-intro}
Assume \eqref{uni-eq:GL}.  If $p>1$ and $f\geq0$ belongs to
$L^p(X,\omega^n)$ with positive integral, there is exactly one normalized
pair $(c,u)$, with $c>0$, $u\in\PSH(X,\omega)\cap L^\infty$,
$\sup_Xu=0$, satisfying
\[
  (\omega+\ddc u)^n=c f\,\omega^n.
\]
Moreover
\[
   c=\frac{\int_X\omega^n}{\int f\,\omega^n},
\]
and $u$ is continuous.  Thus two bounded solutions with the same
$L^p$ density differ by a constant; no positive lower bound on $f$ is
required.
\end{theorem}

We first establish the comparison principle through its equivalence with mass invariance.

\begin{theorem}\label{uni-thm:GL-equivalence}
The following conditions on the same equation form $\omega$ are
equivalent:
\begin{enumerate}
\item $\ddc\omega=0$ and $\ddc(\omega^2)=0$;
\item $\ddc\omega=0$ and $\mathrm{d}\omega\wedge \dc \omega=0$;
\item $\ddc(\omega^k)=0$ for $1\leq k\leq n-1$;
\item $\int_X\mathrm{MA}_\omega(a)=\int_X\omega^n$ for every bounded
      $\omega$-plurisubharmonic $a$;
\item for every bounded $\omega$-plurisubharmonic pair $a,b$,
\begin{equation}\label{uni-eq:ordinaryCP}
 \int_{\{a<b\}}\mathrm{MA}_\omega(b)
 \leq\int_{\{a<b\}}\mathrm{MA}_\omega(a).
\end{equation}
\end{enumerate}
The same equivalence holds when the tests in the last two conditions are
restricted to smooth potentials.  For $n=1$ all conditions are automatic;
for $n=2$ they reduce to $\ddc\omega=0$.
\end{theorem}

\begin{proof}
The differentiation formula
\begin{equation}\label{uni-eq:power-torsion}
 \ddc(\omega^k)=k\omega^{k-1}\wedge\ddc\omega
 +k(k-1)\omega^{k-2}\wedge \mathrm{d}\omega\wedge \dc \omega
\end{equation}
proves the equivalence of the first three conditions.  In dimensions one
and two the omitted equations vanish by degree, giving the stated
caveats.  The characterization of smooth mass invariance is the polarization argument of \cref{thm:mass-invariance-characterization}. Exact-background smooth decreasing regularization and bounded product continuity extend the identity to every bounded potential.

Here is a direct proof of bounded comparison from that identity. Set $q=\max(a,b)$. On $\{a<b\}$, locality gives $\mathrm{MA}_\omega(q)=\mathrm{MA}_\omega(b)$. On $\{a\geq b\}$, one has $a=q$ and $a\leq q$, so ordered contact (\cref{prop:ordered-contact-general}) gives
\[
 \mathbf1_{\{a\geq b\}}\mathrm{MA}_\omega(a)
 \leq \mathbf1_{\{a\geq b\}}\mathrm{MA}_\omega(q).
\]
All three total masses equal $V=\int_X\omega^n$. Subtracting the last inequality from this finite common mass proves \eqref{uni-eq:ordinaryCP}.

Conversely, bounded comparison implies bounded mass invariance: shift one of a bounded pair by a sufficiently large constant so
that its strict sublevel set is all of $X$, apply
\eqref{uni-eq:ordinaryCP}, and reverse the pair.  Comparing with the zero
potential yields the mass identity.  The bounded mass identity restricted
to smooth potentials implies the differential conditions by \cref{thm:mass-invariance-characterization}, closing all implications.
\end{proof}

For the uniqueness argument we need an auxiliary solution and the mixed inequality. Continuous existence supplies the first. If $G\geq0$ lies in $L^p(X,\omega^n)$ and $\int_XG\omega^n=V$, \cref{thm:kn-continuous} gives a continuous bounded solution with a positive normalizing constant. The mass identity makes that constant equal to one. Thus
\begin{equation}\label{uni-eq:aux-solvability}
 \mathrm{MA}_\omega(\rho)=G\omega^n
\end{equation}
is solvable, including for the indicator-truncated density below.

The bounded mixed inequality, proved from local Bedford--Taylor theory in \cref{lem:nonclosed-mixed-density}, applies to bounded $\omega$-plurisubharmonic functions $a,b$, nonnegative integrable densities $F,G$ and $0<t<1$. If $\mathrm{MA}_\omega(a)\geq F\omega^n$ and $\mathrm{MA}_\omega(b)\geq G\omega^n$, it gives
\begin{equation}\label{uni-eq:mixed-bounded}
 \mathrm{MA}_\omega((1-t)a+tb)
 \geq\bigl((1-t)F^{1/n}+tG^{1/n}\bigr)^n\omega^n,
 \qquad 0<t<1.
\end{equation}

\begin{proof}[of Theorem~\ref{uni-thm:GL-intro}]
The domination statement needed here is the case $c=0$ of \cref{thm:bounded-strict-domination}: if $\mathrm{MA}_\omega(a)(\{a<b\})=0$ for bounded $\omega$-plurisubharmonic functions $a,b$, then $a\geq b$.

Let $\mathrm{MA}_\omega(u)=\mathrm{MA}_\omega(v)=\mu=F\omega^n$, with
$F\in L^p$, $p>1$.  The mass identity gives $\mu(X)=V$.
If $u-v$ is nonconstant, shift $u$ so that its difference from $v$
has supremum $s>0$ and infimum $-s$.  Interchanging the pair if
necessary, arrange that $S=\{v<u\}$ satisfies $\mu(S)\leq V/2$.
Define
\[
 b_0=\frac{V-2\mu(S)}V\geq0,
 \qquad G=2F\mathbf1_S+b_0.
\]
This nonnegative $L^p$ function has total integral $V$, so choose the
continuous bounded solution $\rho$ of \eqref{uni-eq:aux-solvability}.
Take $C_0$ with $\rho-C_0\leq u$, and set
\[
 q_\delta=(1-\delta)u+\delta(\rho-C_0),\qquad
 E_\delta=\{v<q_\delta\}\subset S,
 \qquad 0<\delta<1.
\]
On $S$, the mixed inequality gives
\[
 \mathrm{MA}_\omega(q_\delta)\geq
 \bigl((1-\delta)+\delta2^{1/n}\bigr)^n\mu=:A_\delta\mu,
 \qquad A_\delta>1.
\]
Ordinary comparison on $E_\delta$ yields
\[
 A_\delta\mu(E_\delta)
 \leq\int_{E_\delta}\mathrm{MA}_\omega(q_\delta)
 \leq\int_{E_\delta}\mathrm{MA}_\omega(v)=\mu(E_\delta).
\]
Thus $\mu(E_\delta)=0$.  Bounded domination gives
$v\geq q_\delta$ everywhere; letting $\delta\downarrow0$ gives
$v\geq u$, contradicting $\sup_X(u-v)=s>0$.  Therefore the
difference is constant.  The argument never divides by $F$ and permits it to vanish on an open set.

Finally, for the density $f$ in the theorem put
$I=\int_Xf\omega^n>0$ and $c=V/I$.  Equation
\eqref{uni-eq:aux-solvability} supplies a continuous bounded solution for
$G=cf$, normalized by subtracting its supremum.  Any competing bounded
solution with constant $c'$ satisfies $V=c'I$, hence $c'=c$;
the just-proved uniqueness and the common normalization identify the
potentials.  In particular every bounded solution with this $L^p$
density is continuous, even though continuity was not assumed.
\end{proof}

\begin{remark}
Theorem~\ref{uni-thm:GL-equivalence} is an equivalence of differential,
comparison, and mass-invariance properties.  It is not an equivalence
between the Guan--Li condition and uniqueness.  The other criteria in this chapter apply without that background condition.  Conversely, imposing it
on an auxiliary metric rather than on $\omega$ does not justify
\eqref{uni-eq:ordinaryCP} for $\mathrm{MA}_\omega$.
\end{remark}
